\section{Which Telemetry Identifies the Target?}
\subsection{Inspecting the field at its producer}
We examined two sets of physical demonstration logs to ask whether a replay program could identify the target accepted by the receiver. In 12 DROID recordings~\cite{droid2024}, the observation-time and action-time values of \texttt{prev\_command\_successful} differed on 81 of 2,610 non-skipped rows. A second set of 13 recordings from disjoint site--date pairs had 176 differences in 4,551 rows, including 105 in one recording. All 25 recordings stored calculated Cartesian commands. None included the Franka native commanded-pose field or rolling reception rate. In the archived Polymetis client, a change in desired torque sets the status flag.\footnote{\href{https://github.com/facebookresearch/fairo/blob/0a01a7fa7a7c65b2f9a3aebf5e79040940daf9d2/polymetis/polymetis/src/clients/franka_panda_client/franka_panda_client.cpp}{Polymetis Franka client, archived implementation.}} Reconstructing a target therefore requires the producer, timestamp and meaning of each field.

\subsection{Reading the application cycle of a software filter}
To isolate the information provided by application time, we compiled the translational equilibrium filter from the Franka impedance example into a local cross-process harness.\footnote{\href{https://github.com/frankarobotics/franka_ros/blob/35e1f654426e04bc9f83b73af4ab68a3fb145c84/franka_example_controllers/src/cartesian_impedance_example_controller.cpp}{Franka Cartesian impedance example, pinned implementation.}} The client and controller communicated over local UDP. The controller processed callbacks every cycle or every ten cycles, with three repetitions per setting and 900 updates per run. Each observer used only messages received by the decision time. Across 5,400 updates, the application-cycle acknowledgments for all 120 commands arrived before the first affected target sample, allowing reconstruction to numerical precision. Across the three repetitions, using send time to infer application time gave target RMS errors of 0.074--0.144\,mm with callbacks every cycle and 0.585--0.791\,mm with callbacks every ten cycles (Fig.~\ref{fig:origin}b). Median round-trip times for requested targets ranged from 0.802 to 1.487\,ms across six runs. Adding the private equilibrium to 27 status messages per run cost 1,675--1,712 payload bytes with zero extra requests. Both the field and the acknowledgments were supplied by our harness.

\subsection{When does a bit identify the applied input?}
An execution acknowledgment can substitute for a target value if it identifies the applied displacement. On 24 fresh Sawyer layouts we sent either three execution bits or the stored target every three steps, with both packets arriving two steps later. For binary forwarding or substitution by zero, knowledge of the initial target, applied increments and workspace clip was sufficient. Both streams produced identical trajectories in all 96 paired conditions. When five accepted translations were subsequently attenuated to one quarter, the bits continued to report forwarding. Mean hand RMS then reached 10.126\,mm with execution bits, 5.149\,mm with targets and 6.679\,mm with empty packets. Each policy still completed all 48 attenuation trials. In 5,025 of 7,200 decisions under attenuation, interpreting a forwarded command as full execution caused the predicted origin box to exclude the true origin. These are correlated control steps from 24 layouts.
